\scanpage{112}
% archive foliation: 111
le fibré $\mathcal T/B_u$ \inserted{sur $P$} ? C'est le fibré \uline{principal associé} à $\Omega^1_{P/X}$ \inserted{(ou encore des ``éléments de contact projectifs d'ordre 1'')}. Une section en dehors de $\sigma$ est donc \deleted{une \uline{trivialisation} de $\underline\omega$} un couple d'un quotient inversible $L$ de $P^1_{X/S}(\underline\omega)$ (qui donne $\sigma'$) -- i.e. un op. diff.
\[
D:\underline\omega\unsure{\text{---}|\longrightarrow}L\deleted{\unsure{inversible}}\inserted{\text{d'ordre 1}}
\]
\uline{engendrant} $L$, \uline{plus} une trivialisation $\varphi$ de $\deleted{\unsure{\underline\omega}}\inserted{N\otimes L^{-1}}$, où $N:=\operatorname{Ker}(P^1_{X/S}(\underline\omega)\longrightarrow L))$.
\margin{$N\otimes L\simeq\det P^1_{X/S}(\underline\omega)\simeq\underline\omega^3$ donc $N\otimes L^{-1}\simeq L^{-2}\otimes\underline\omega^3$ d'où $\varphi:L^2\xrightarrow{\sim}\underline\omega^3$. i.e. $L$ est une ``racine carrée'' de $\underline\omega^3$.}
\deleted{\unsure{op. diff.} $D:\underline\omega\longrightarrow\mathcal O_X$ engendrant $\mathcal O_X$}

Le cas $\sigma'$ disjoint de $\sigma$ correspond aux couples $(D,\unsure{\lambda|\varphi})$ d'une \uline{connexion} de $\underline\omega$, $D:\underline\omega\longrightarrow\underline\omega^2$, \uline{et} d'une trivialisation $\varphi$ de $\unsure{\underline\omega^2|\underline\omega}$, $\varphi\in\Gamma\underline\omega$.
Condition d'horizontalité ? \deleted{Il faut que} cette trivialisation soit horizontale ! \unsure{pour la connexion $D$ sur $\underline\omega$} \deleted{\unsure{comme ci-dessus}}, i.e. $D\varphi=0$.
\deleted{\unsure{Ramener à une condition exprimée sur $\underline\omega$}}

En termes de la connexion projective $c_0$, quel est le faisceau $\mathcal U\subset\operatorname{Applhol}(X^{\mathrm{an}},\mathbb C^{\mathrm{an}})$ ? C'est le faisceau des \inserted{germes} d'applications \inserted{étales} $t:X^{\mathrm{an}}\longrightarrow\mathbb P^1_{\mathbb C}(\mathbb C)$ telles que $t^*(c)=c_0$, où $c$ est la conn. projective canonique (et unique !) sur $\mathbb P^1(\mathbb C)$.
\margin{NB $t^*(c)=c_{D_t}$, et \unsure{par suite} on doit $c_{D_t}=c_0$.}

\deleted{\unsure{Si}} Si $s:X'^{\mathrm{an}}\longrightarrow\mathbb E^1(\mathbb C)\subset\mathbb P^1(\mathbb C)$, est étale on peut considérer la connexion \inserted{sur $\underline\omega$} $s^*(D)\deleted{\unsure{\underline\omega}}$, où $D$ est la conn. canonique sur $\Omega^1_{\mathbb P^1(\mathbb C)/\mathbb C}$ \deleted{$\unsure{\varphi}$} (la connexion $D_T$, où $T$ est la coordonnée \unsure{canonique}) c'est aussi $D_s$. On peut alors, \deleted{\unsure{pour}} pour une connexion \deleted{\unsure{régulière}} $D:\underline\omega\longrightarrow\underline\omega^2$ donnée \deleted{telle que $c_D=c_0$}, \uline{regarder} les $t$ telles que $D=D_s$. \deleted{\unsure{Cela}} Soit $t$ un uniformisant, \unsure{alors on a donc}
\[
\underline D=D_t+\bar\omega\quad(\bar\omega\in\Gamma(\underline\omega)),
\]
et l'équation en $s$ est $D_s=D_t+\bar\omega$, or on sait que
\[
D_s=D_t+\deleted{\unsure{ds}}\underbrace{s'\,d\frac1{s'}}_{-\frac{ds'}{\unsure{s'|s'^2}}},\qquad s':=\frac{ds}{dt}.
\]
Alors l'équation est
\[
-\frac{ds'}{s'}=\bar\omega.
\]
